\chapter{Gestion des Risques}

\section{Identification et Prévention}

La gestion des risques est une composante essentielle de la gestion de projet. Nous avons anticipé plusieurs risques :

\begin{table}[H]
\centering
\begin{tabularx}{\textwidth}{L{4cm} C{2cm} C{2cm} X}
\toprule
\textbf{\color{primarycolor}Risque} &
\textbf{\color{primarycolor}Probabilité} &
\textbf{\color{primarycolor}Impact} &
\textbf{\color{primarycolor}Action Préventive} \\
\midrule
Perte de données (Kafka tombe) & Moyenne & Élevé & Configuration de la rétention et mécanisme de Haute Disponibilité (réplication). \\
\addlinespace[0.15cm]
Retard de livraison & Élevée & Élevé & Utilisation du Planning Poker pour des estimations réalistes. \\
\addlinespace[0.15cm]
Incohérence des résultats & Faible & Élevé & Mise en place du test de validation de cohérence (Replay vs Temps réel). \\
\bottomrule
\end{tabularx}
\caption{Matrice des risques et actions préventives}
\end{table}

\section{Le Cône d'Incertitude}

Au début du projet, l'incertitude liée à la technologie Kafka était très forte. Au fur et à mesure des Sprints, la maîtrise technique s'est améliorée et l'incertitude s'est réduite, permettant des estimations plus précises pour les tâches restantes.

\begin{warningbox}
    \textbf{\textcolor{red!70!black}{\large Risque Technique}}\\[0.3cm]
    Les technologies Big Data/Streaming sont complexes à configurer. L'utilisation de \texttt{docker-compose} a permis d'encapsuler cette complexité et d'atténuer le risque de problèmes d'environnement.
\end{warningbox}
